DNA methylation is a stable, information-rich epigenetic signal that reflects tissue identity
and serves as a powerful biomarker for disease diagnosis and biological age assessment.
However, many widely used analytical approaches rely on linear models, which cannot directly
capture nonlinear dependencies and context-dependent interactions across the methylome. Here
we introduce \method{}, a Transformer-based foundation model pretrained without phenotype
supervision on 169{,}120 profiles from diverse tissues and studies, represented at 49{,}156
CpGs. To learn global rather than locally constrained representations, we adapt Whole-Cell
Expression Decoding to methylation as the Whole-CpG Encoder--Decoder (WCED). WCED encodes
two independently sampled partial views of each profile and requires their CLS representations
to reconstruct CpGs absent from the corresponding inputs, while symmetric InfoNCE aligns views
of the same sample and separates different samples. On held-out profiles, \method{} retained
sample identity across partial measurements, reconstructed unseen CpGs substantially better than
population and representation-disruption controls, and captured age- and tissue-associated
variation without labels. Fine-tuning improved chronological-age prediction over matched neural
and linear baselines. These findings establish \method{} as a reusable DNA methylation
foundation model while revealing study-specific structure as a central challenge for cross-cohort
generalization.

% Internal drafting note (not printed): expand demonstrated downstream breadth before submission.
